\subsection*{Estudio técnico}
\addcontentsline{toc}{subsection}{Estudio técnico}

\subsubsection*{Tamaño del Proyecto y Capacidad Operativa}
\addcontentsline{toc}{subsubsection}{Tamaño}

La determinación del tamaño de la organización se fundamenta en un análisis dual: el cumplimiento del marco legal colombiano y la capacidad técnica necesaria para el despliegue de la plataforma. De acuerdo con el \textit{Decreto 957 de 2019}, la empresa se clasifica inicialmente como una \textbf{Microempresa}, dado que sus ingresos proyectados en el sector de Servicios se sitúan por debajo del límite de 32.988 UVT anuales \parencite{Decreto1297}.

Desde la perspectiva de la ingeniería de sistemas, el tamaño se define por la arquitectura requerida para garantizar la disponibilidad y el procesamiento de datos adaptativos. Para asegurar un desempeño eficiente, se consideran los siguientes factores:

\begin{itemize}
	\item \textbf{Recursos Humanos Especializados:} La estructura prioriza un departamento de desarrollo capaz de gestionar el ciclo de vida del software y el entrenamiento de modelos de \textit{Machine Learning}, fundamentados en las tendencias de automatización de asesoría financiera \parencite{BBVAResearch01}.
	
	\item \textbf{Potencial de Cómputo:} Al ser una solución \textit{EdTech} \parencite{UniversidadEuropea01}, el tamaño técnico se dimensiona para una concurrencia controlada, utilizando infraestructura en la nube que permite un escalamiento horizontal según la demanda del mercado de nómadas digitales y jóvenes en Bogotá.
	
	\item \textbf{Factores de Producción Digital:} La materia prima consiste en activos digitales y hardware de alto rendimiento necesarios para que los motores de recomendación operen con baja latencia, garantizando la eficiencia técnica validada en pruebas de usuario.
\end{itemize}

A continuación, se presenta la tabla de clasificación empresarial vigente en Colombia. Para el año 2026, el valor de la Unidad de Valor Tributario (UVT) se establece en \$ 52.374 pesos colombianos.

\begin{adjustbox}{
		center,
		caption=[{Clasificación de empresas en Colombia.}]{\centering Clasificación de empresas en Colombia. \textit{Fuente:} Decreto 957 de 2019, Gobierno Nacional de Colombia.},
		label={ClasificacionEmpresasColombia},
		nofloat=table, vspace={20px}}
	\resizebox{\textwidth}{!}{
		\begin{tabular}{|c|c|c|c|}
			\hline
			\rowcolor[HTML]{D9EAD3} 
			\textbf{Tamaño} & \textbf{Manufactura} & \textbf{Servicios} & \textbf{Comercio} \\ \hline
			\textbf{Microempresa} & \begin{tabular}[c]{@{}c@{}}Hasta 23.563 \\ UVT\end{tabular} & \begin{tabular}[c]{@{}c@{}}Hasta 32.988 \\ UVT\end{tabular} & \begin{tabular}[c]{@{}c@{}}Hasta 44.769 \\ UVT\end{tabular} \\ \hline
			\begin{tabular}[c]{@{}c@{}} \textbf{Pequeña} \\ \textbf{Empresa}\end{tabular} & \begin{tabular}[c]{@{}c@{}}Desde 23.563 UVT \\ hasta 204.995 UVT\end{tabular} & \begin{tabular}[c]{@{}c@{}}Desde 32.988 UVT \\ hasta 131.951 UVT\end{tabular} & \begin{tabular}[c]{@{}c@{}}Desde 44.769 UVT \\ hasta 431.196 UVT\end{tabular} \\ \hline
			\begin{tabular}[c]{@{}c@{}} \textbf{Mediana} \\ \textbf{Empresa}\end{tabular} & \begin{tabular}[c]{@{}c@{}}Desde 204.995 UVT \\ hasta 1'736.565 UVT\end{tabular} & \begin{tabular}[c]{@{}c@{}}Desde 131.951 UVT \\ hasta 483.034 UVT\end{tabular} & \begin{tabular}[c]{@{}c@{}}Desde 431.196 UVT \\ hasta 2'160.692 UVT\end{tabular} \\ \hline
		\end{tabular}
	}
\end{adjustbox}

La adopción de este marco regulatorio permite a la organización alinearse con las políticas nacionales de fomento empresarial y desarrollo sostenible. Dado que el núcleo del negocio es una plataforma virtual dedicada a la alfabetización financiera de jóvenes, independientes y nómadas digitales, la empresa se ubica estrictamente en el sector de Servicios. Por su etapa temprana y enfoque geográfico inicial en Bogotá, se ratifica su categorización como microempresa, con una hoja de ruta técnica diseñada para escalar hacia niveles superiores de productividad y alcance tecnológico.

\subsubsection*{Localización}
\addcontentsline{toc}{subsubsection}{Localización}

La localización estratégica es un factor determinante para asegurar el posicionamiento y la escalabilidad del servicio. Aunque la plataforma opera en un entorno digital, su centro de operaciones se establece en la ciudad de Bogotá D.C., decisión respaldada por la alta penetración de tecnologías de la información y comunicación en la capital \parencite{DANE_Internet2023}. Esta ubicación permite conectar directamente con el segmento de jóvenes y trabajadores informales que presentan las mayores brechas de alfabetización financiera en la región \parencite{DANE_Juventud2025}.

La sede operativa se selecciona bajo criterios de eficiencia técnica y administrativa:

\begin{itemize}
	\item \textbf{Conectividad e Infraestructura:} Acceso a redes de alta velocidad y servicios básicos estables, indispensables para el mantenimiento de servidores y la gestión de bases de datos financieras.
	
	\item \textbf{Acceso al Talento Humano:} Proximidad a centros de innovación y universidades, facilitando la colaboración en el desarrollo de algoritmos de personalización y diseño de experiencia de usuario (UX).
	
	\item \textbf{Ecosistema Fintech:} Bogotá se posiciona como el hub principal para el desarrollo de tecnologías financieras en Colombia \parencite{Finnovista2024}, lo que facilita la creación de alianzas estratégicas con entidades bancarias y reguladores locales.
\end{itemize}

\paragraph*{Lugares candidatos}
\addcontentsline{toc}{subsubsection}{Lugares candidatos}

Para dar con la elección de los lugares candidatos, se tienen en cuenta los espacios diseñados para trabajar colaborativamente dentro de ambientes agradables y que colaboren con el flujo de las actividades continuas y sostenidas, actualmente espacios como los \textbf{coworking} logran cumplir aquella necesidad y la infraestructura buscada, además de contar junto a espacios en los que se pueda manejar la interdisciplinariedad de las distintas partes que comparten el ambiente de trabajo, lo que permite aportar valor al producto y mejorar las oportunidades de apoyo de nuestro modelo. Los lugares candidatos son:

\begin{itemize}
	\item \textbf{WeWork:} Son espacios de trabajo compartidos cuyo objetivo es ofrecer soluciones de trabajo en las que todos podamos ser más productivos, motivados y lograr generar mejores proyectos. Además de contar con soluciones de espacio para pequeñas y medianas empresas, inclusive a grandes empresas también. Los cuales no solamente se limita a una oficina en particular sino que también logran ofrecer sus servicios a distintas ciudades de Latinoamérica.
	
	\begin{adjustbox}{
			center,
			caption=[{Oficinas de WeWork.}]{\centering Oficinas de WeWork. \textit{Fuente:} WeWork.com.co},
			label={TablaWeWork},
			nofloat=table, vspace={20px}}
			\begin{tabular}{|c|c|}
				\hline
				\rowcolor[HTML]{D9EAD3} 
				\textbf{Servicio} & \textbf{Precio x Mes} \\ \hline
				Plan All Access & \$390.000 + IVA \\ \hline
				Plan All Access Plus & \$600.000 + IVA \\ \hline
			\end{tabular}
	\end{adjustbox}
	
	\item \textbf{Office ToGo:} Son espacios de trabajo compartidos dentro de los cuales las pequeñas y medianas empresas y emprendedores y trabajadores independientes tienen la oportunidad de poder utilizar espacios de trabajo equipados y organizados con todo lo necesario, además de encontrarse en las mejores ubicaciones de Bogotá.
	
	\begin{adjustbox}{
			center,
			caption=[{Oficinas de OfficeToGo.}]{\centering Oficinas de OfficeToGo. \textit{Fuente:} OfficeToGo.co},
			label={TablaOfficeToGo},
			nofloat=table, vspace={20px}}
			\begin{tabular}{|c|c|}
				\hline
				\rowcolor[HTML]{D9EAD3} 
				\textbf{Servicio} & \textbf{Precio x Mes} \\ \hline
				Planes virtuales & \$99.000 + IVA \\ \hline
				Plan emprendedor & \$169.000 + IVA \\ \hline
				Oficinas generales & \$1.100.000 + IVA \\ \hline
				Oficina semi-privada & \$850.000 + IVA \\ \hline
			\end{tabular}
	\end{adjustbox}
	
	\item \textbf{ParqueSoft:} ParqueSoft es un proveedor multisectorial de servicios de las TIC (tecnologías de la información) que cuenta con el apoyo de más de 50 empresas de base dentro de Bogotá, que trabajando juntas generan una oferta de valor integral logrando un ecosistema de apoyo, buscando crear soluciones de calidad, a la vanguardia de las últimas tendencias digitales. Además ParqueSoft es el clúster de ciencia y tecnología informática más grande de Colombia y uno de los líderes más importantes de temas de apoyo a proyectos de emprendimiento con base tecnológica, impulsándolos y colaborando a que crezcan.
	
	\begin{adjustbox}{
			center,
			caption=[{Oficinas de ParqueSoft.}]{\centering Oficinas de ParqueSoft. \textit{Fuente:} ParqueSoftBogota.com},
			label={TablaParqueSoft},
			nofloat=table, vspace={20px}}
			\begin{tabular}{|c|c|}
				\hline
				\rowcolor[HTML]{D9EAD3} 
				\textbf{Servicio} & \textbf{Precio x Mes} \\ \hline
				Pago Mensual & \$96.000 al mes $\sim$ \$1.152.000 al año \\ \hline
				Pago Semestral & \$78.000 al mes $\sim$ \$932.000 al año \\ \hline
				Pago Anual & \$69.000 al mes $\sim$ \$828.000 al año \\ \hline
			\end{tabular}
	\end{adjustbox}

	\item \textbf{Elección Final:} Tras analizar las variables de costo, ubicación, servicios técnicos y valor agregado, se selecciona a \textbf{ParqueSoft Bogotá} como la sede inicial del proyecto. Esta decisión se fundamenta en la alineación de ParqueSoft con la naturaleza tecnológica de la plataforma. La posibilidad de acceder a capacitaciones especializadas, mentorías en ingeniería de software y un ecosistema de innovación resulta crítica para la maduración del Producto Mínimo Viable (PMV) y la implementación de los modelos de Machine Learning.
	
\end{itemize}

\subsubsection*{Tipo de emprendimiento}
\addcontentsline{toc}{subsubsection}{Tipo de emprendimiento}

El proyecto se clasifica bajo el modelo de Startup, adoptando la metodología \textit{Lean Startup}. Este enfoque se centra en el desarrollo ágil y en la iteración constante del producto a partir del aprendizaje validado con el usuario final. Al tratarse de una solución tecnológica emergente que integra inteligencia artificial, esta metodología permite mitigar el riesgo de mercado mediante un diseño centrado en el cliente, optimizando el uso de recursos técnicos y financieros, y facilitando una arquitectura escalable desde sus etapas iniciales.

La implementación de este marco de trabajo asegura que la plataforma evolucione en respuesta directa a las necesidades del entorno financiero de Bogotá, transformando la retroalimentación en requisitos funcionales de software de manera acelerada.

\begin{figure}[H]
	\centering
	\includegraphics[width=12cm]{Imagenes/MetodologiaLeanStartUp.png}
	\caption[{Metodología Lean Startup.}]{\centering Metodología Lean Startup. \textit{Fuente:} Adaptado de ThePowerMBA (2025).}
	\label{fig:metodologialeanstartup}
\end{figure}

Bajo este ciclo iterativo de ``Aprender-Medir-Construir``, se propone el desarrollo de un \textit{Producto Mínimo Viable (PMV)}. Este prototipo funcional permite someter a prueba las hipótesis de valor y crecimiento con una inversión mínima, evaluando la aceptabilidad técnica de los simuladores y motores de recomendación. En caso de que las métricas de usabilidad o el modelo de negocio no alcancen los umbrales esperados, la metodología facilita la adaptación o ajuste estratégico del enfoque tecnológico sin comprometer la viabilidad global del emprendimiento.

\subsubsection*{Equipo de trabajo}
\addcontentsline{toc}{subsubsection}{Equipo de trabajo}

Dentro de la planeación e implementación de este proyecto, estamos conformados por:

\begin{itemize}
	\item \textbf{Gerente general:} El encargado de planificar las actividades que ocurran dentro de la empresa. Organizar los recursos dentro de la entidad y definir en que dirección se va a dirigir la empresa en el corto, mediano y largo plazo, entre muchas otras tareas más.
	
	\item \textbf{Líder del área de ingeniería:} Responsable de liderar un equipo de desarrollo y el responsable acerca de la calidad de sus producto. Busca establecer una visión técnica con el equipo de desarrollo y trabaja con ellos para buscar alcanzar el objetivo.
	
	\item \textbf{Director del departamento de investigación, desarrollo e innovación (I+D+I):} Es la persona encargada de dirigir la investigación, la innovación y el desarrollo. Por lo tanto es responsable del avance tecnológico e investigativo centrados hacia el avance de la sociedad, siendo una de las partes más elementales dentro de las tecnologías informativas. Junto a este rol, también tiene la tarea de coordinar y manejar los equipos interdisciplinares para desarrollar nuevos productos o tecnologías.
\end{itemize}

A continuación el equipo de trabajo será el siguiente:

\begin{adjustbox}{
		center,
		caption=[{Equipo de trabajo.}]{\centering Equipo de trabajo \textit{Fuente:} Autores},
		label={TablaEquipoTrabajo},
		nofloat=table, vspace={20px}}
	\resizebox{\textwidth}{!}{
		\begin{tabular}{|c|c|}
			\hline
			\rowcolor[HTML]{D9EAD3} 
			\textbf{Integrante} & \textbf{Rol} \\ \hline
			David Leonardo Rincón Ovalle & Gerente general, líder del área de ingeniería \\ \hline
			Leidy Marcela Morales Segura & Director del departamento de investigación, desarrollo e innovación \\ \hline
		\end{tabular}
	}
\end{adjustbox}

